% Gerado por scripts/migrate_markdown_to_latex.py.
% Fonte migrada: docs/apresentacoes/simposio_2026/roteiros/roteiro_10_minutos.md


\chapter{Roteiro de 10 Minutos}



Tempo total planejado: aproximadamente 9min40s.



\section{Slide 1 - Capa (20s)}


Objetivo: abrir a fala e posicionar o tema.






Fala: Este trabalho discute eficiência computacional em operações de Transformers. O foco está em quantização, pruning e em como comparar essas estratégias de forma experimental. Hoje eu apresento o problema, a hipótese e o desenho do estudo.




Transição: Primeiro, vale situar por que esse tema é relevante.



\section{Slide 2 - Roteiro (20s)}


Objetivo: dar previsibilidade para a banca.







Fala: Vou dividir a apresentação em cinco partes. Começo pelo contexto e pelo problema. Depois passo pelas operações centrais e pelas estratégias de otimização. Em seguida mostro a metodologia, o estado atual do projeto e as conclusões.




Transição: Começando pelo contexto geral.



\section{Slide 3 - Transformers e a IA atual (35s)}


Objetivo: justificar relevância.






Fala: Transformers se tornaram a arquitetura dominante em linguagem, visão e multimodalidade. Isso acontece porque eles combinam paralelização com boa modelagem de dependências de longo alcance. Com essa adoção ampla, eficiência deixou de ser detalhe técnico e passou a ser requisito prático.




Transição: Mas o problema aparece de forma mais clara quando olhamos a inferência.



\section{Slide 4 - Custo da inferência (30s)}


Objetivo: mostrar por que otimizar é necessário.






Fala: O custo da inferência cresce com o comprimento da sequência, com a dimensão interna do modelo e com o número de camadas. Em produção, isso aparece como latência, uso de memória e limite de throughput. Então, tornar a inferência mais eficiente é essencial para viabilizar esses modelos em hardware real.




Transição: Esse custo se concentra em alguns gargalos bem definidos.



\section{Slide 5 - Gargalos principais (40s)}


Objetivo: localizar os gargalos.






Fala: Os gargalos centrais são a autoatenção, as projeções densas e o KV cache. A atenção cresce com o contexto, as projeções densas repetem multiplicações matriciais e o KV cache pressiona a memória. Esses pontos são justamente os mais promissores para otimização.




Transição: Isso leva ao problema de pesquisa do trabalho.



\section{Slide 6 - Problema e hipótese (30s)}


Objetivo: explicitar o recorte científico.






Fala: O problema de pesquisa é reduzir custo computacional sem degradar excessivamente a saída. O recorte escolhido foi trabalhar com projeções densas e self-attention em cenários controlados de inferência. A hipótese inicial é que quantização tende a oferecer um trade-off melhor do que pruning.




Transição: Com esse recorte, os objetivos ficam mais claros.



\section{Slide 7 - Objetivos (25s)}


Objetivo: resumir o que será feito.






Fala: O trabalho implementa um protótipo simplificado das operações centrais do Transformer. Depois compara baseline, pruning por magnitude e quantização INT8 simulada. Por fim, mede custo computacional e qualidade de saída.




Transição: Agora eu mostro as operações estudadas.



\section{Slide 8 - Projeções densas (30s)}


Objetivo: explicar a primeira operação.






Fala: As projeções densas transformam vetores por multiplicações com matrizes de pesos. Elas aparecem na construção de Q, K e V e também em camadas lineares do bloco Transformer. Por isso, já concentram custo importante mesmo antes de analisar a atenção completa.




Transição: A segunda operação é a self-attention com KV cache.



\section{Slide 9 - Self-attention e KV cache (40s)}


Objetivo: explicar a segunda operação.






Fala: Na self-attention, o modelo gera Q, K e V, calcula escores, aplica softmax e monta a saída contextualizada. Para inferência autoregressiva, o KV cache evita recomputação, mas cresce com o tamanho do contexto. Então essa operação impacta tanto latência quanto memória.




Transição: Com essas operações definidas, entram as estratégias de otimização.



\section{Slide 10 - Quantização (30s)}


Objetivo: apresentar a primeira estratégia.






Fala: Quantização reduz a precisão numérica usada para representar pesos, ativações ou caches. Os ganhos esperados são menor uso de memória e maior throughput. O ponto de atenção é o erro numérico introduzido na saída.




Transição: A outra estratégia relevante é o pruning.



\section{Slide 11 - Pruning (30s)}


Objetivo: contrastar com quantização.






Fala: Pruning remove pesos considerados menos relevantes e introduz esparsidade. A economia teórica pode ser alta, mas o ganho real depende muito do suporte do hardware. Além disso, o efeito na saída tende a ser menos previsível.




Transição: A literatura recente reforça por que a hipótese inicial favorece quantização.



\section{Slide 12 - TurboQuant (45s)}


Objetivo: sustentar a hipótese com literatura.







Fala: Um trabalho importante aqui é o TurboQuant. Ele foi divulgado pelo Google Research em 24 de março de 2026, com paper disponível no arXiv desde 28 de abril de 2025. O foco está em compressão de vetores e do KV cache para inferência. Para o meu projeto, isso não é prova final, mas reforça a decisão de priorizar quantização no primeiro ciclo experimental.




Transição: Com essa base teórica, eu passo para a metodologia.



\section{Slide 13 - Metodologia: desenho (35s)}


Objetivo: mostrar o desenho do experimento.






Fala: O estudo compara três cenários: baseline, pruning em 50\% e quantização INT8 simulada. As operações analisadas são projeção densa e self-attention de cabeça única. A grade principal varia comprimento de sequência e dimensão para observar o comportamento do custo.




Transição: Além do desenho, também defini métricas e controles.



\section{Slide 14 - Metodologia: métricas (35s)}


Objetivo: mostrar rigor experimental.






Fala: O ambiente previsto usa Python 3.11, bibliotecas numéricas e GPU RTX 4070 Ti Super. Cada cenário terá aquecimento, repetições e sincronização para medir latência. As métricas cobrem latência, memória, FLOPs e qualidade de saída com MSE, MAE, R² e cosseno.




Transição: Com isso, eu mostro o que já está pronto no projeto.



\section{Slide 15 - Resultados parciais (30s)}


Objetivo: mostrar andamento real.






Fala: O projeto ainda não tem benchmarks completos executados. Mas o código-base já implementa operações, métricas e cenários. Ou seja, o trabalho já tem estrutura técnica pronta para iniciar a coleta comparativa.




Transição: Mesmo antes dos dados próprios, a literatura já aponta uma direção inicial.



\section{Slide 16 - Evidências da literatura (40s)}


Objetivo: usar o gráfico como motivação.






Fala: Esse gráfico resume evidências publicadas sobre TurboQuant. De um lado aparece a redução relativa de memória para KV cache; do outro, o ganho relativo de velocidade reportado. Esses números não são resultados do meu projeto, mas ajudam a justificar por que quantização é a hipótese inicial mais forte.




Transição: Com isso, eu posso fechar as conclusões.



\section{Slide 17 - Conclusões (35s)}


Objetivo: fechar com direção clara.






Fala: O trabalho ataca um problema real de inferência em Transformers. Hoje, a quantização aparece como a estratégia inicial mais promissora para esse recorte. O próximo passo é validar essa hipótese com benchmarks próprios e análise comparativa dos cenários.




Transição: Por fim, deixo as referências principais.



\section{Slide 18 - Referências (20s)}


Objetivo: mostrar base bibliográfica.





Fala: Aqui estão as referências principais que sustentam a apresentação. Elas incluem o trabalho original de Transformer, revisões recentes de compressão e as fontes usadas para o slide de TurboQuant.




Transição: Encerrando, fico à disposição para perguntas.



\section{Slide 19 - Obrigado (10s)}


Objetivo: fechamento.





Fala: Obrigado pela atenção. Fico aberto para perguntas e discussão.



\section{Onde acelerar se faltar tempo}


\begin{itemize}
\item Resumir o Slide 4 em uma frase.
\item Juntar verbalmente os Slides 8 e 9.
\item Resumir o Slide 18 em uma única frase.
\end{itemize}
